\section{对唯象学的意义}
\label{sec:pheno}

本节给出 NNPDF3.1 PDF 唯象学意义的初步研究。
%
首先，我们在上一节讨论的基础上总结 PDF 不确定度的状况；讨论 PDF 不确定度之后，
我们聚焦于奇异味与粲 PDF。接着讨论作为强子对撞机过程首要输入的 PDF 亮度，
以及 LHC 过程的预言，特别是 LHC 的 $W$、$Z$ 与希格斯产生。
与其他各处一样，这里只展示部分结果，更大规模的图表集可在线获取（见第~\ref{sec:delivery} 节）。

\subsection{PDF 不确定度的改进}

\label{sec:phenopdfs}

在第~\ref{sec:impactnewdata} 节逐一讨论了每条新数据对 NNPDF3.1 PDF 的影响并分离出新方法学的效应之后，
现在可以通过比较 NNPDF3.0 与 NNPDF3.1 的 PDF 不确定度来研究全部新数据的综合效应。
图~\ref{fig:pdfunc31vs30} 进行了这一比较：图中比较 NNPDF3.0 与
NNPDF3.1 集合 PDF 的相对不确定度（全部以共同归一化计算），
表示为 up 与 down 的价夸克（即 $q-\bar q$）与海（即 $\bar q$），以及 strange 与 charm 的 $q+\bar q$。
结果既对单个 PDF 味道给出，也对文献~\cite{Ball:2016neh} 式 (3.4) 定义的 singlet、valence 与 triplet 组合给出。

最显眼的效应是胶子不确定度的大幅缩小：如今在几乎所有 $x$ 处都达到百分之一量级。
如第~\ref{sec:impactnewdata} 节所讨论，这源于来自 DIS（尤其在 HERA）、$Z$ 横向动量分布、喷注产生与顶夸克对产生的许多相互一致的约束的共同作用——合起来覆盖了非常宽的运动学范围。
与胶子混合的单态夸克组合在所有 $x\lsim 0.1$ 显示出相当的改进，但在大 $x$ 较不明显。

有意思的是，对夸克 PDF 而言，味基与文献~\cite{Ball:2016neh} 式 (3.4) 给出的“演化”基下的不确定度模式并不相同。具体而言，
前述单态组合不确定度的缩小并未出现在任何轻夸克价分布上——它们在 NNPDF3.1 与 NNPDF3.0 中不确定度大体相当。
这是因为由于粲如今独立参数化，NNPDF3.1 的味分离多少更不确定。
更多的实验信息（特别是 LHCb 与 ATLAS 数据）弥补了这一点，但在小 $x$ 与极大 $x$ 处则不然。
实际上除 $x\sim0.1$ 外，valence 与 triplet 分布在 NNPDF3.1 中的不确定度总体略大于 NNPDF3.0。

%%%%%%%%%%%%%%%%%%%%%%%%%%%%%%%%%%%%%%%%%%%%%%%%%%%%%%%%%%%%%%%%%%%%%
\begin{figure}[t!]
\begin{center}
  \includegraphics[scale=0.38]{images/plots/valence_nnpdf30nnlo.pdf}
  \includegraphics[scale=0.38]{images/plots/valence_nnpdf31nnlo.pdf}
  \includegraphics[scale=0.38]{images/plots/sea_nnpdf30nnlo.pdf}
  \includegraphics[scale=0.38]{images/plots/sea_nnpdf31nnlo.pdf}
  \includegraphics[scale=0.38]{images/plots/evolbasis3_nnpdf30nnlo.pdf}
  \includegraphics[scale=0.38]{images/plots/evolbasis3_nnpdf31nnlo.pdf}
  \caption{\small 
NNPDF3.0（left）与 NNPDF3.1
  （right) NNLO PDF 相对不确定度的比较，以 NNPDF3.1
  NNLO 中心值归一。top 为两个轻夸克价 PDF 与 gluon；center 为全部海 PDF；
   bottom 为 singlet、valence 与同位旋 triplet 组合。
    %
\label{fig:pdfunc31vs30}}
\end{center}
\end{figure}
%%%%%%%%%%%%%%%%%%%%%%%%%%%%%%%%%%%%%%%%%%%%%%%%%%%%%%%%%%%%%%%%%%%%%%%

这一模式也见于海夸克不确定度。在小 $x\lsim 10^{-2}$ 它们彼此相当，
且都（奇异味除外）比 NNPDF3.0 对应值小得多，因为它们受小 $x$ 微扰演化中单态与胶子混合的驱动~\cite{das}。
如今独立参数化的 charm PDF 也具有更小的不确定度。
在大 $x$ 区则轮到较不精确的味分离知识发挥作用，相对不确定度更大。
显然在 NNPDF3.0 中
charm 曾有不自然的小不确定度，因为大 $x$ 的微扰生成粲被绑在胶子上。而在 NNPDF3.1 中，
价区内海 PDF 不确定度的层级符合预期：up 与 down 最精确，
strange 与 charm 则承受越来越大的不确定度。
NNPDF3.1 中 up 尤其是 down 海分量的不确定度略有增大，但与 NNPDF3.0 仍相当；
strangeness 的不确定度保持稳定，charm 的则显著增大——鉴于大 $x$ 粲基本不受数据约束，这本应如此。
就此而言，一个有趣的现象是：如上文第~\ref{sec:emc} 节所讨论，通过纳入 EMC 数据集可以在大 $x$
获得对粲及其他海 PDF 更准确的确定。未来关于 $Z+c$ 产生等过程的 LHC 数据或可检验 EMC 数据集的可靠性。

\subsection{奇异味 PDF}
\label{sec:strangeness}

up 与 down PDF 的大小（即中心值）存在广泛共识——现有确定在其小不确定度内相互符合良好；
而奇异味 PDF 的大小却一直存在争议，这里结合 NNPDF3.1 结果重新审视。
具体地，质子夸克海的奇异味占比定义为
\be
\label{eq:rs}
R_s (x,Q^2)=\frac{ s(x,Q^2)+\bar{s}(x,Q^2) }{ \bar{u}(x,Q^2)+\bar{d}(x,Q^2) } \, ,
\ee
相应的动量分数比定义为
\be
\label{eq:rsint}
K_s (Q^2)=\frac{\int_0^1 dx\, x \lp s(x,Q^2)+\bar{s}(x,Q^2)\rp }{ 
\int_0^1 dx\, x\lp \bar{u}(x,Q^2) + \bar{d}(x,Q^2)\rp }  \, ,
\ee
传统上被认为明显小于一；在尚不能从数据提取奇异味 PDF 的年代产生的 PDF 集合（例如 NNPDF1.0~\cite{Ball:2008by}）
常对所有 $x$ 取 $R_s\sim\frac{1}{2}$，从而也有 $K_s\sim\frac{1}{2}$。这种程度的奇异味压低确实见于许多近期全局 PDF 集合中——其中对奇异味的最大约束来自深度非弹性中微子单举 $F_2$ 与粲
$F_2^c$（“双 $\mu$”）数据。

文献~\cite{Aad:2012sb}对这一图景提出挑战：基于 ATLAS $W$ 与 $Z$ 产生数据并结合 HERA DIS
数据，该文主张在测量区域内奇异味占比 $R_s$ 为一左右的量级。分别基于 NNPDF2.3 与 NNPDF3.0 全局分析的
文献~\cite{Ball:2012cx,Ball:2015oha}（两者都包含文献~\cite{Aad:2012sb}的数据）得出结论：
ATLAS 数据虽然确实偏好更大的奇异味 PDF，但由于不确定度很大，其对全局 PDF 确定的影响是温和的；
而且若只用 HERA 与 ATLAS 数据确定奇异味 PDF，中心值与文献~\cite{Aad:2012sb}的结论一致，
但不确定度大到足以在一倍标准差内与全局 PDF 集合的压低奇异味相符。
文献~\cite{Ball:2015oha}还表明 CMS $W+c$ 产生数据~\cite{Chatrchyan:2013uja}——那里首次纳入且也包含在 NNPDF3.1 中（但因缺乏 NNLO 修正的知识而只在 NLO 拟合中）——由于不确定度大而影响可忽略。

如第~\ref{sec:atlaswz} 节所讨论，ATLAS $W$ 与 $Z$ 产生数据已由精确得多的
文献~\cite{Aaboud:2016btc}数据集补充，后者同样据称偏好增强的奇异味。的确，
第~\ref{sec:atlaswz} 节已看到纳入这些数据显著增强了奇异味；
第~\ref{sec:results-mc} 节还看到，得益于 NNPDF3.1 新特征——独立参数化的粲 PDF，
这一增强可以被全局 PDF 确定所容纳。
因此重新评估 NNPDF3.1 中的奇异味很有意义：我们用理论上动机明确的不同数据集选择进行比较：
与先前用默认 NNPDF3.1 得到的 NNPDF3.0 奇异味结果相比，
我们比较第~\ref{sec:collideronly} 节仅对撞机集合（理论上看更可靠），以及我们构造的一个 PDF 集合——
采用 NNPDF3.1 方法但只纳入表~\ref{tab:completedataset} 全部 HERA 单举结构函数数据与文献~\cite{Aaboud:2016btc}
的 ATLAS 数据。由于单举 DIS 数据本身无法单独确定奇异味~\cite{Forte:2010dt}，
这是一个完全依赖 ATLAS 数据的奇异味确定。

表~\ref{tab:rs} 给出用这些不同 PDF 集合得到的 NNLO 结果：
$Q=1.38~{\rm GeV}$（低于粲阈值）与 $Q=m_Z$、$x=0.023$ 处的 $R_s(x,Q)$ 式~\ref{eq:rs}——
$x=0.023$ 是 ATLAS 为最大化敏感性而选的值。结果也与文献~\cite{Aaboud:2016btc}比较。
表的图形化展示见图~\ref{fig:xRsplot}。

%%%%%%%%%%%%%%%%%%%%%%%%%%%%%%%%%%%%%%%%%%%%%%%%%%%%%%%%%%%%%%%%%%%
\begin{table}[t]
  \begin{center}
    \small
        \renewcommand{\arraystretch}{1.2}
\begin{tabular}{|l|c|c|}
\hline
\centering PDF set &  $R_s(0.023,1.38~\mathrm{GeV})$ & $R_s(0.023,M_{Z})$ \\
\hline
\hline
NNPDF3.0    & 0.45$\pm$0.09    & 0.71$\pm$0.04 \\
NNPDF3.1     & 0.59$\pm$0.12  & 0.77$\pm$0.05 \\
NNPDF3.1 collider-only   & 0.82$\pm$0.18  & 0.92$\pm$0.09 \\
NNPDF3.1 HERA + ATLAS $W,Z$   & 1.03$\pm$0.38  & 1.05$\pm$0.240\\	
\hline
{\tt xFitter} HERA + ATLAS $W,Z$ (Ref.~\cite{Aaboud:2016btc}) 	& $1.13
\,{}^{+0.11}_{-0.11}$& - \\
\hline
\end{tabular}
\end{center}
\caption{\small \label{tab:rs} 
  低标度与高标度下 $x=0.023$ 处的奇异味占比 $R_s(x,Q)$ 式~(\ref{eq:rs})。
  图中给出 NNPDF3.0 以及 NNPDF3.1 基线、仅对撞机与 HERA+ATLAS $W,Z$
  集合的结果，并与 {\tt xFitter} ATLAS 数值（文献~\cite{Aaboud:2016btc}）比较。
}
\end{table}
%%%%%%%%%%%%%%%%%%%%%%%%%%%%%%%%%%%%%%%%%%%%%%%%%%%%%%%%%%%%%%%%%%%

%%%%%%%%%%%%%%%%%%%%%%%%%%%%%%%%%%%%%%%%%%%%%%%%%%%%%%%%%%%%%%%%%%%%%
\begin{figure}[t]
\begin{center}
  \includegraphics[scale=0.35]{images/plots/Rsplot.pdf}
  \includegraphics[scale=0.35]{images/plots/Rsplot-highQ.pdf}
  \caption{\small 表~\ref{tab:rs} 结果的图形化展示。
\label{fig:xRsplot}}
\end{center}
\end{figure}
%%%%%%%%%%%%%%%%%%%%%%%%%%%%%%%%%%%%%%%%%%%%%%%%%%%%%%%%%%%%%%%%%%%%%%%

首先，把 NNPDF3.1 HERA+ATLAS $W,Z$ 结果与基于相同数据的文献~\cite{Aaboud:2016btc}比较可见一倍标准差水平的一致性，
中心值相近而不确定度大约增至四倍——最可能是因为更灵活的参数化以及独立参数化粲。
其次，NNPDF3.1 的奇异味明显大于 NNPDF3.0：如第~\ref{sec:results-mc}、\ref{sec:atlaswz} 节所示，
这主要归因于 ATLAS~$W,Z$~2011 数据的作用加上由数据确定粲；
从图~\ref{fig:31-nnlo-old-vs-new} 清楚可见新数据与新方法学都导致奇异味增强，前者主导而后者不可忽略。
这种增强在仅对撞机集合中更为突出，其数值非常接近来自 ATLAS 数据的值。
这提示对撞机数据偏好的奇异味与数据集其余部分（最可能是中微子数据）之间存在某种张力。


把这一分析推广到整个 $x$ 范围是有意义的。
图~\ref{fig:xRs-abs-atlaswz2011} 完成：$R_s(x,Q)$ 式~(\ref{eq:rs})
按低高两种标度随 $x$ 作图，此时只包含 NNPDF3.0 以及 NNPDF3.1 的默认与仅对撞机版本。
显然在仅对撞机 PDF 集合中大 $x$ 的奇异味几乎无约束，
而全局拟合受中微子数据约束呈压低值 $R_s\sim 0.5$。更低 $x$ 处可见它与对撞机数据约束之间的张力——后者偏好更大的值。

%%%%%%%%%%%%%%%%%%%%%%%%%%%%%%%%%%%%%%%%%%%%%%%%%%%%%%%%%%%%%%%%%%%%%
\begin{figure}[t]
\begin{center}
  \includegraphics[scale=0.37]{images/plots/xRs-abs-atlaswz2011lowQ.pdf}
  \includegraphics[scale=0.37]{images/plots/xRs-abs-atlaswz2011highQ.pdf}
   \includegraphics[scale=0.37]{images/plots/xRs-abs-atlaswz2011lowQ-global.pdf}
  \includegraphics[scale=0.37]{images/plots/xRs-abs-atlaswz2011highQ-global.pdf}
  \caption{\small 奇异味比 $R_s(x,Q)$ 式~(\ref{eq:rs})
    随 $x$ 的变化，两个标度 $Q=1.38$ GeV（left）
    与 $Q=m_Z$（right)。
    %
     top 比较 NNPDF3.1 与仅对撞机 NNPDF3.1；bottom 与 CT14 及 MMHT 比较。
\label{fig:xRs-abs-atlaswz2011}}
\end{center}
\end{figure}
%%%%%%%%%%%%%%%%%%%%%%%%%%%%%%%%%%%%%%%%%%%%%%%%%%%%%%%%%%%%%%%%%%%%%%%

图~\ref{fig:xRs-abs-atlaswz2011} 还将 NNPDF3.1 的奇异味比 $R_s(x,Q)$
与 CT14 及 MMHT14 比较。
%
我们发现整个 $x$ 范围内一致性良好，
而 NNPDF3.1 的 PDF 误差通常小于另两家，尤其在高标度处。
同样有趣的是 NNPDF3.1 中比值 $R_s$ 的 PDF 不确定度在极大 $x$ 处急剧膨胀，
反映出该运动学区域缺乏关于奇异味的直接信息。


%%%%%%%%%%%%%%%%%%%%%%%%%%%%%%%%%%%%%%%%%%%%%%%%%%%%%%%%%%%%%%
\begin{table}[t]
  \begin{center}
        \renewcommand{\arraystretch}{1.2}
\begin{tabular}{|l|c|c|}
\hline
\centering PDF set &  $K_s(Q=1.38~\mathrm{GeV})$ & $K_s(Q=M_{\rm Z})$ \\
\hline
\hline
NNPDF3.0   &  $0.45 \pm 0.07$  & $0.72 \pm 0.04$ \\
NNPDF3.1    & $0.53 \pm 0.07$  &  $ 0.75 \pm 0.04 $  \\
NNPDF3.1 collider-only   & $3.4 \pm 2.5$  & $1.5 \pm 0.6$  \\
NNPDF3.1 HERA + ATLAS $W,Z$   & $-1.0 \pm 7.0$  & $2.8 \pm 1.7$ \\
\hline
\end{tabular}
\end{center}
\caption{\small \label{tab:rsint} 
  低标度与高标度下的奇异味动量分数式~(\ref{eq:rsint})。
  图中给出 NNPDF3.0 以及 NNPDF3.1 基线、仅对撞机与 HERA+ATLAS $W,Z$ PDF 集合的结果。}
\end{table}
%%%%%%%%%%%%%%%%%%%%%%%%%%%%%%%%%%%%%%%%%%%%%%%%%%%%%%%%%%%%%%%%%


现在转向奇异味动量分数 $K_s(Q^2)$ 式~(\ref{eq:rsint})；相同 PDF 集合与标度的数值列于表~\ref{tab:rsint}。结果与表~\ref{tab:rs} 分析所得相当类似。
对仅对撞机、尤其是 HERA + ATLAS $W,Z$ 拟合，$K_s$ 的
中心值非物理且不确定度巨大；基本上只能说奇异味动量分数 $K_s$ 完全不确定。
这相当戏剧性地表明表~\ref{tab:rs} 中相对精确的数值只在一个相当窄的 $x$ 范围内成立。
未来的更多 LHC 数据能否带来有竞争力的仅对撞机拟合，从而确认奇异味增强并在更宽的 $x$ 范围内精确测定奇异味，值得期待。

\subsection{再论质子的粲含量}
\label{sec:phenocharm}

通过拟合粲 PDF 确定的质子粲含量已在 NNPDF 全局分析框架内于文献~\cite{Ball:2016neh}中量化，那里首次比较了粲独立参数化与微扰生成两种结果。
%
该分析只在 NLO 进行，数据集与 NNPDF3.0 拟合非常相近。
%
我们现在在微扰 QCD 的 NNLO 层面、并在纳入新数据集——特别是顶夸克以及对粲 PDF 有显著影响和约束的 LHCb、ATLAS 电弱玻色子产生数据——的背景下重新审视拟合粲 PDF。

事实上第~\ref{sec:disentangling} 节（特别是图~\ref{fig:31-nnlo-old-vs-new}）已经看到：
NNPDF3.1 新增的数据显著缩小了粲 PDF 不确定度，但也影响其中心值——在大 $x$ 变化超过一个标准差。
另外，文献~\cite{Ball:2016neh}中大 $x$ 粲主要受 EMC 数据约束（见第~\ref{sec:emc} 节讨论），
而这些数据未进入默认 NNPDF3.1 PDF 确定。如第~\ref{sec:emc} 节所见这些数据对粲仍有重要影响，
因此在包含与不包含该数据集两种情形下都应对粲确定进行重新评估。
于是现在比较三种结果：默认 NNPDF3.1 NNLO 集合、
第~\ref{sec:results-mc} 节所述微扰生成粲的修改版，以及第~\ref{sec:emc} 节所述在 NNPDF3.1
数据集中追加 EMC 数据的修改版。

表~\ref{tab:momsr} 给出粲动量分数
\be
\label{eq:charmmomfrac}
C(Q^2) \equiv \int_0^1\, dx\,\lp xc(x,Q^2)+x\bar{c}(x,Q^2)\rp \ ,
\ee
在两个 $Q^2$ 值——粲阈值处 $Q=m_c$ 与 $Q=m_Z$——用这些 PDF 集合算得的数值。
表~\ref{tab:momsr} 结果的图形化展示见图~\ref{fig:xCplot}。
根据粲是独立参数化还是微扰生成，动量分数的不确定度相差两个数量级。
不过当粲被参数化时，中心值在巨大的不确定度内彼此一致，
对应中心值只是略大于微扰粲情形（尽管按微扰生成结果的不确定度尺度衡量差异巨大）。
加入 EMC 数据后不确定度缩小约三倍，中心值略有增大，与该数据集对粲 PDF 的影响（第~\ref{sec:emc} 节）一致。

我们可以把“粲独立参数化并由数据确定”与“粲微扰生成”两种总动量分数之差解释为
“内禀”（即非微扰）粲动量分数。
当粲参数化并计入 EMC 数据时，在 $Q=m_c$ 我们发现它为
$C(m_c)_{\rm FC}-C(m_c)_{\rm PC} =  \left(0.16 \pm 0.14\right)$\%：这为一倍标准差水平上质子内小的内禀粲分量提供了证据，
一定程度上改进了文献~\cite{Ball:2016neh}的估计，但仍带有相当大的不确定性。
%
我们对质子非微扰粲分量的估计明显小于 CT14IC 模型分析~\cite{Dulat:2013hea}允许的范围（亦见图~\ref{fig:xCplot}）。
但它大于文献~\cite{Jimenez-Delgado:2014zga}的预期——那里声称四倍标准差水平的上限为 $0.5$\%。
这两项分析都难以拟合 EMC 粲结构函数数据，原因在于其粲 PDF 函数形式过于受限。
在高标度，所有 $C(Q)$ 的估计都被小 $x$ 粲 PDF 的微扰增长主导，
但含 EMC 数据拟合中的一倍标准差超额仍然存在，尽管占整体的份额不断减小。
这一点在图~\ref{fig:momfrac} 中得到非常清楚的展示，图中画出粲动量分数 $C(Q^2)$ 对标度 $Q$ 的依赖。


%%%%%%%%%%%%%%%%%%%%%%%%%%%%%%%%%%%%%%%%%%%%%%%%%%%%%%%%%%%%%%%%%%%%%%%%
\begin{table}[t]
  \center
   \renewcommand{\arraystretch}{1.4}
  \begin{tabular}{|l|c|c|}
    \hline
    PDF set & $C(m_c)$ & $C(m_Z)$ \\
    \hline
    \hline
    NNPDF3.1           & $\lp 0.26 \pm 0.42\rp $\%  &  $\lp 3.8 \pm 0.3\rp $\%\\
    NNPDF3.1 with pert. charm        & $\lp 0.176 \pm 0.005\rp $\%  &  $\lp 3.73 \pm 0.02\rp $\% \\
    NNPDF3.1 with EMC data  &   $\lp 0.34 \pm 0.14\rp $\% & $\lp 3.8  \pm 0.1\rp $\% \\
          \hline
  \end{tabular}
  \caption{\small \label{tab:momsr}
    恰在粲阈值上方 $Q=m_c$~GeV 与 $Q=m_Z$ 处的粲动量分数 $C(Q^2)$ 式~(\ref{eq:charmmomfrac})。
    结果针对 NNPDF3.1 及其两个修改版：向数据集加入 EMC 数据者，或不拟合粲者。
    %
  }
  \end{table}
%%%%%%%%%%%%%%%%%%%%%%%%%%%%%%%%%%%%%%%%%%%%%%%%%%%%%%%%%%%%%%%%%%%%%%%%

%%%%%%%%%%%%%%%%%%%%%%%%%%%%%%%%%%%%%%%%%%%%%%%%%%%%%%%%%%%%%%%%%%%%%
\begin{figure}[t]
\begin{center}
  \includegraphics[scale=0.35]{images/plots/Cplot.pdf}
  \includegraphics[scale=0.35]{images/plots/Cplot-highQ.pdf}
  \caption{\small 表~\ref{tab:momsr} 中 $C(Q^2)$
    结果在 $Q=m_c$ GeV（left）与 $Q=m_Z$（right) 的图形化展示。
    文献~\cite{Dulat:2013hea}的模型估计亦给出。
\label{fig:xCplot}}
\end{center}
\end{figure}
%%%%%%%%%%%%%%%%%%%%%%%%%%%%%%%%%%%%%%%%%%%%%%%%%%%%%%%%%%%%%%%%%%%%%%%



%%%%%%%%%%%%%%%%%%%%%%%%%%%%%%%%%%%%%%%%%%%%%%%%%%%%%%%%%%%%%%%%%%%%%
\begin{figure}[t]
\begin{center}
  \includegraphics[scale=0.30]{images/plots/momfrac.pdf}
  \caption{\small 表~\ref{tab:momsr} 的粲动量分数对标度 $Q$ 作图。
\label{fig:momfrac}}
\end{center}
\end{figure}
%%%%%%%%%%%%%%%%%%%%%%%%%%%%%%%%%%%%%%%%%%%%%%%%%%%%%%%%%%%%%%%%%%%%%%%

这些粲动量分数取值的起源可以通过直接比较粲 PDF 来理解（图~\ref{fig:PDFlargexcharm}），
再次取恰在阈值上方 $Q=m_c$ 与 $Q=m_Z$ 两处，后者以对基线 NNPDF3.1 结果之比表示。
粲动量分数在微扰生成与参数化由数据确定两种情形下的一致性，与如下事实有关：参数化时的最优拟合粲
具有与 NNLO 微扰生成粲定性相同的形状——正如文献~\cite{Ball:2016neh}及上文第~\ref{sec:results-mc}
节所观察到的。
然而在阈值 $Q=m_c$ GeV 处，最优拟合粲在大 $x$（$x\gsim 0.2$）大于微扰分量，
虽带有较大不确定度——加入 EMC 数据集后有所减小而形状无显著变化。
加入 EMC 数据后，在中小 $10^{-2}\lsim x\lsim 10^{-1}$ 区域 PDF 被推向加入前的误差带上沿，
不确定度大幅缩小。微扰生成粲 PDF 小得不切实际的不确定度显而易见；
EMC 数据使 $x\gsim 10^{-3}$ 的不确定度缩小也清晰可见——这在第~\ref{sec:emc} 节已讨论过。

%%%%%%%%%%%%%%%%%%%%%%%%%%%%%%%%%%%%%%%%%%%%%%%%%%%%%%%%%%%%%%%%%%%%%
\begin{figure}[t]
  \begin{center}
    \includegraphics[scale=0.35]{images/plots/xc-abs-charmcomp-lowQ.pdf}
  \includegraphics[scale=0.35]{images/plots/xc-charmcomp-highQ.pdf}
  \includegraphics[scale=0.35]{images/plots/xc-ERR-charmcomp-lowQ.pdf}
  \includegraphics[scale=0.35]{images/plots/xc-ERR-charmcomp-highQ.pdf}
    \caption{\small 表~\ref{tab:momsr} 标度与 PDF 集合下 charm PDF 的比较。
      top 为 PDF 本身，bottom 为相对不确定度。
\label{fig:PDFlargexcharm}}
\end{center}
\end{figure}
%%%%%%%%%%%%%%%%%%%%%%%%%%%%%%%%%%%%%%%%%%%%%%%%%%%%%%%%%%%%%%%%%%%%%%%

粲 PDF 在微扰生成与参数化由数据确定两种情形之间的差异来源，以及把后者分解为“内禀”
（非微扰）与微扰分量的可能方案，可以通过研究阈值附近对标的依赖来理解——类比文献~\cite{Ball:2016neh}中的类似分析。
图~\ref{fig:charmPDFnf4} 完成了这项工作：图中给出（$n_f=4$ 方案下）粲 PDF 随 $x$ 的变化，
粲分别为参数化与微扰生成。
一方面，$x\gsim0.1$ 的大 $x$
 粲 PDF 分量本质上与标度无关：微扰粲在该区域严格为零，因此该区域的拟合结果
 可解释为具有非微扰起源，即“内禀”。

另一方面对更小 $x$ 粲 PDF 强烈依赖标度。微扰生成时它在 $Q\sim 2$~GeV 已对所有 $x\gsim 10^{-3}$
相当可观且为正，而在阈值处对一切 $x\lsim10^{-2}$ 变得又大又负——可能因为大的未被重求和的小 $x$ 贡献。
最优拟合参数化粲 PDF 在其较大不确定度内，除微扰粲变号的那个依赖标度的点之外，
本质上对所有 $x\lsim 10^{-1}$ 都更平坦且模更小。
对所有 $x\lsim 0.1$，拟合粲与微扰粲之间这种形状差异远大于二者各自的不确定度。
这一观察似乎支持如下结论：粲为微扰生成的假设可能是一种偏差来源。

%%%%%%%%%%%%%%%%%%%%%%%%%%%%%%%%%%%%%%%%%%%%%%%%%%%%%%%%%%%%%%%%%%%%%
\begin{figure}[t]
\begin{center}
  \includegraphics[scale=0.37]{images/plots/charm-nf4-Qdep-log.pdf}
  \includegraphics[scale=0.37]{images/plots/charm-nf4-Qdep-lin.pdf}
  \includegraphics[scale=0.37]{images/plots/charm-nf4-Qdep-log-nFC.pdf}
  \includegraphics[scale=0.37]{images/plots/charm-nf4-Qdep-lin-nFC.pdf}
  \caption{\small $n_f=4$ 方案下小 $x$
    （left）与大 $x$（right plot）的 charm PDF 在不同 $Q$ 值下的行为：
    NNPDF3.1 NNLO PDF 集合（top）与假设粲微扰生成（bottom）。
\label{fig:charmPDFnf4}}
\end{center}
\end{figure}
%%%%%%%%%%%%%%%%%%%%%%%%%%%%%%%%%%%%%%%%%%%%%%%%%%%%%%%%%%%%%%%%%%%%%%%

更新后的 NNPDF3.1 分析与我们早期粲研究的结论一致~\cite{Ball:2016neh}，即
质子的非微扰粲分量与全局数据相容，只是我们现在纳入的新高精度对撞机数据对该分量可能的大小给出了更严格的界限。
%
特别是一旦加入 EMC 粲结构函数数据集，从表~\ref{tab:momsr} 可见约 $\simeq 0.5\%$
的非微扰粲动量分数已相当于偏离最优拟合值二到三个标准差。
因此携带质子动量高达 1\% 的内禀粲模型已被现有数据强烈排斥。

\subsection{部分子亮度}
\label{sec:lumis}

在对各个 PDF 及其不确定度的唯象意义作解析讨论之后，现在转向驱动强子对撞机过程的
部分子亮度（定义依文献~\cite{Mangano:2016jyj}）。
图~\ref{fig:lumi-31-vs-30} 比较 NNPDF3.0 与 NNPDF3.1 NNLO 集合在 LHC~13 TeV 的部分子亮度，
图~\ref{fig:lumi-nnpdf31-2d} 以二维等高线图展示其不确定度随末态不变质量 $M_y$
与快度 $y$ 的分布，全部以 NNPDF3.1
中心值归一。我们给出 quark-quark、quark-antiquark、gluon-gluon 与 gluon-quark
亮度的结果，它们对应不与个别味道耦合的末态（如 $Z$ 或希格斯）的测量。
不确定度图中还给出 up-antidown 亮度作参照（与例如 $W^+$
产生相关）。

%%%%%%%%%%%%%%%%%%%%%%%%%%%%%%%%%%%%%%%%%%%%%%%%%%%%%%%%%%%%%%%%%%%%%
\begin{figure}[t]
\begin{center}
  \includegraphics[scale=0.38]{images/plots/q2_31-vs-30.pdf}
  \includegraphics[scale=0.38]{images/plots/qq_31-vs-30.pdf}
  \includegraphics[scale=0.38]{images/plots/gg_31-vs-30.pdf}
   \includegraphics[scale=0.38]{images/plots/qg_31-vs-30.pdf}
   \caption{\small NNPDF3.0 与 NNPDF3.1 NNLO PDF 集合在 LHC~$13$ TeV 的部分子亮度比较。
     自左至右、自上而下依次为 quark-antiquark、quark-quark、gluon-gluon
     与 quark-gluon 亮度。
     结果以 NNPDF3.1 中心值归一给出。
\label{fig:lumi-31-vs-30}}
\end{center}
\end{figure}
%%%%%%%%%%%%%%%%%%%%%%%%%%%%%%%%%%%%%%%%%%%%%%%%%%%%%%%%%%%%%%%%%%%%%%%


%%%%%%%%%%%%%%%%%%%%%%%%%%%%%%%%%%%%%%%%%%%%%%%%%%%%%%%%%%%%%%%%%%%%%
\begin{figure}[h]
\begin{center}
   \includegraphics[scale=0.40]{images/plots/qqbar_plot_lumi2d_uncertainty_30.pdf}   
   \includegraphics[scale=0.40]{images/plots/qqbar_plot_lumi2d_uncertainty.pdf}   
\includegraphics[scale=0.40]{images/plots/qq_plot_lumi2d_uncertainty_30.pdf}
\includegraphics[scale=0.40]{images/plots/qq_plot_lumi2d_uncertainty.pdf}
\includegraphics[scale=0.40]{images/plots/gg_plot_lumi2d_uncertainty_30.pdf}
\includegraphics[scale=0.40]{images/plots/gg_plot_lumi2d_uncertainty.pdf}
  \includegraphics[scale=0.40]{images/plots/gq_plot_lumi2d_uncertainty_30.pdf}
  \includegraphics[scale=0.40]{images/plots/gq_plot_lumi2d_uncertainty.pdf}
\includegraphics[scale=0.40]{images/plots/udbar_plot_lumi2d_uncertainty_30.pdf}
\includegraphics[scale=0.40]{images/plots/udbar_plot_lumi2d_uncertainty_31.pdf}
   \caption{\small 图~\ref{fig:lumi-31-vs-30} 各亮度的相对不确定度，
     表示为末态不变质量 $M_X$ 与快度 $y$ 的函数；left 列为 NNPDF3.0，right 列为 NNPDF3.1
     （上四排）。bottom row 给出 up-antidown 亮度的结果。
\label{fig:lumi-nnpdf31-2d}}
\end{center}
\end{figure}
%%%%%%%%%%%%%%%%%%%%%%%%%%%%%%%%%%%%%%%%%%%%%%%%%%%%%%%%%%%%%%%%%%%%%%%


这一比较有两个明显特征。第一，对所有不变质量夸克亮度普遍变大，
而相比 NNPDF3.0，NNPDF3.1 的胶子亮度在较小不变质量处略有增强、较大不变质量处略有压低。
夸克部门位移的大小约为一倍标准差有时甚至更大，而胶子一般明显低于一倍标准差。
这当然反映了第~\ref{sec:PDFcomparisons} 节所见 PDF 的模式，
特别见图~\ref{fig:31-nnlo-vs30}。
第二，相比 NNPDF3.0，NNPDF3.1 的不确定度大幅缩小。
图~\ref{fig:lumi-nnpdf31-2d} 中这一缩减令人印象深刻：清楚可见 NNPDF3.0 时大部分相空间的不确定度量级为
5\%，而现在在宽的中心快度范围
$|y|\lsim 2$ 内、对 $100$~GeV$\lsim M_x\lsim 1$~TeV 的末态质量约为 1--2\%。
这是上文第~\ref{sec:phenopdfs} 节讨论的胶子与夸克单态 PDF 不确定度缩小的直接后果。
确实，对味分解敏感的亮度——如图~\ref{fig:lumi-nnpdf31-2d} 同样给出的 up-antidown 亮度——
从 NNPDF3.0 到 NNPDF3.1 并未表现出显著的不确定度缩小。


接下来我们把 NNPDF3.1 与 CT14、MMHT14 比较：结果示于图~\ref{fig:lumi-31-vs-global}。
对 quark-quark 亮度我们发现符合良好；quark-antiquark 中心值的离散稍大但仍在
一倍标准差水平相合。
在大质量 $M_X\gsim 2$ TeV 处符合变得勉强，反映大 $x$ PDF 知识有限。
对 gluon-gluon 与 gluon-quark 道，在与 LHC 精密物理相关的 $M_X\simeq 600$ GeV 以下符合合理，
而在与 BSM 搜寻相关的更大质量处——尤其 NNPDF3.1 与 MMHT14 之间——符合较差。
%
当然应当记住：与 MMHT14 和 CT14 相比，NNPDF3.1 数据集更宽、独立参数化 PDF 更多，
一旦所有全局 PDF 集合都完成更新，形势可能改变。

接着在图~\ref{fig:lumi-31-vs-abmp} 与 ABMP16 PDF 比较。此处同时给出对应 ABMP16
默认集合（$\alpha_s(m_Z)=0.1147$）与迄今所有比较采用的共同 $\alpha_s(m_Z)=0.118$ 的结果。
    %
    当使用 ABMP16 默认值 $\alpha_s(m_Z)=0.1147$ 时 NNPDF3.1 与 ABMP16 存在可观差异（gluon-gluon
    亮度尤甚）；ABMP16 也改用 $\alpha_s(m_Z)=0.118$ 后一致性改善。
    然而 ABMP16 亮度在低与高 $M_X$ 具有极小的不确定度，
   这大概是过度受限的参数化、以及使用不设容忍度的 Hessian 方法的后果（见第~\ref{sec:PDFcomparisons} 节讨论）。

\clearpage

%%%%%%%%%%%%%%%%%%%%%%%%%%%%%%%%%%%%%%%%%%%%%%%%%%%%%%%%%%%%%%%%%%%%%
\begin{figure}[t]
\begin{center}
  \includegraphics[scale=0.38]{images/plots/q2_31-vs-global.pdf}
  \includegraphics[scale=0.38]{images/plots/qq_31-vs-global.pdf}
  \includegraphics[scale=0.38]{images/plots/gg_31-vs-global.pdf}
   \includegraphics[scale=0.38]{images/plots/qg_31-vs-global.pdf}
   \caption{\small 同图~\ref{fig:lumi-31-vs-30}，此处比较 NNPDF3.1 NNLO 与 CT14 及 MMHT14。
\label{fig:lumi-31-vs-global}}
\end{center}
\end{figure}
%%%%%%%%%%%%%%%%%%%%%%%%%%%%%%%%%%%%%%%%%%%%%%%%%%%%%%%%%%%%%%%%%%%%%%%

%%%%%%%%%%%%%%%%%%%%%%%%%%%%%%%%%%%%%%%%%%%%%%%%%%%%%%%%%%%%%%%%%%%%%
\begin{figure}[t]
\begin{center}
  \includegraphics[scale=0.35]{images/plots/q2_31-vs-abmp.pdf}
  \includegraphics[scale=0.35]{images/plots/qq_31-vs-abmp.pdf}
  \includegraphics[scale=0.35]{images/plots/gg_31-vs-abmp.pdf}
   \includegraphics[scale=0.35]{images/plots/qg_31-vs-abmp.pdf}
   \caption{\small 同图~\ref{fig:lumi-31-vs-global}，此处与 ABMP16 PDF 比较，
     分别用其默认 $\alpha_s(m_Z)=0.1149$ 和共同值 $\alpha_s(m_Z)=0.118$。
\label{fig:lumi-31-vs-abmp}}
\end{center}
\end{figure}
%%%%%%%%%%%%%%%%%%%%%%%%%%%%%%%%%%%%%%%%%%%%%%%%%%%%%%%%%%%%%%%%%%%%%%%


%%%%%%%%%%%%%%%%%%%%%%%%%%%%%%%%%%%%%%%%%%%%%%%%%%%%%%%%%%%%%%%%%%%%%
\begin{figure}[t]
\begin{center}
  \includegraphics[scale=0.35]{images/plots/q2_31-fc-vs-pc.pdf}
  \includegraphics[scale=0.35]{images/plots/qq_31-fc-vs-pc.pdf}
  \includegraphics[scale=0.35]{images/plots/gg_31-fc-vs-pc.pdf}
   \includegraphics[scale=0.35]{images/plots/qg_31-fc-vs-pc.pdf}
   \caption{\small 同图~\ref{fig:lumi-31-vs-30}，此处比较 NNPDF3.1
     与其微扰粲修改版。
\label{fig:lumi-31-fc-vs-pc}}
\end{center}
\end{figure}
%%%%%%%%%%%%%%%%%%%%%%%%%%%%%%%%%%%%%%%%%%%%%%%%%%%%%%%%%%%%%%%%%%%%%%%

最后，为进一步强调新 NNPDF3.1 方法学的唯象影响，我们把 NNPDF3.1 PDF
与其微扰生成粲的修改版比较（已在第~\ref{sec:results-mc}、~\ref{sec:phenocharm} 节讨论）。
结果示于图~\ref{fig:lumi-31-fc-vs-pc}。
一方面，我们确认尽管多了一个参数化 PDF，不确定度并未增大。另一方面对中心值的效应温和但不可忽略。
对 gluon-gluon 与 quark-gluon 亮度，差异始终低于一倍标准差且通常更小。
%
对 quark-quark 道，结果在 $M_X\gsim 200$ GeV 不依赖于粲的处理，
但更小不变质量时微扰生成粲给出的 PDF 亮度大于最优拟合参数化粲。
%
对 quark-antiquark 亮度，小 $M_X$ 有类似模式，中等与较大 $M_X$ 处也存在一些差别。




\clearpage

\subsection{LHC 13 TeV 的 $W$ 与 $Z$ 产生}
\label{sec:pheno13TeV}

我们将基于 NNPDF3.1 集合的理论预言与 ATLAS~\cite{Aad:2016naf}在 $\sqrt{s}=13$ TeV 的 $W$ 与 $Z$
产生数据比较。CMS 的类似测量~\cite{CMS-PAS-SMP-15-004}仍属初步故未纳入比较。
%
我们用 {\tt FEWZ}~\cite{Gavin:2012sy}以 NNLO QCD 精度计算 fiducial 截面，分别用
NNPDF3.1、NNPDF3.0、CT14、MMHT14 与 ABMP16 PDF 并附相应 PDF 不确定度带。所有计算（包括 ABMP16）都用 $\alpha_s=0.118$。
      %
      电弱 NLO 修正对 $Z$ 用 {\tt FEWZ} 计算，对 $W$
      用 {\tt HORACE3.2}~\cite{CarloniCalame:2007cd} 计算。
      %
    ATLAS $W^{\pm}$ 截面测量的 fiducial 相空间为：
    带电轻子横向动量与赝快度满足 $p_T^l\ge 25$ GeV 与 $|\eta_l|\le 2.5$，
    缺失能量 $p_T^\nu\ge 25$ GeV，$W$ 横质量 $m_T\ge 50$ GeV。
    %
    对 $Z$ 产生：
    带电轻子横向动量与快度满足 $p_T^l\ge 25$ GeV 与 $|\eta_l|\le 2.5$，双轻子不变质量 $66\le m_{ll} \le 116$ GeV。
          
%%%%%%%%%%%%%%%%%%%%%%%%%%%%%%%%%%%%%%%%%%%%%%%%%%%%%%%%%%%%%%%%%%%%%
\begin{figure}[t]
\begin{center}
  \includegraphics[scale=0.88]{images/plots/WPWM_ratio_ATLAS_comparison_shaded.pdf}
  \includegraphics[scale=0.88]{images/plots/WZ_ratio_ATLAS_comparison_shaded.pdf}
  \caption{\small ATLAS 在 $\sqrt{s}=13$ TeV 测量的 $W^+/W^-$ 比（left）
    与 $W/Z$ 比（right)
    与不同 NNLO PDF 集合理论预言的比较。
    %
    预言按正文所述分别给出含（heavy）与不含（light）用 {\tt FEWZ}、{\tt HORACE}
    计算的 NLO EW 修正两种。
\label{fig:LHC-13tev-incxsec-1}}
\end{center}
\end{figure}
%%%%%%%%%%%%%%%%%%%%%%%%%%%%%%%%%%%%%%%%%%%%%%%%%%%%%%%%%%%%%%%%%%%%%%%

图~\ref{fig:LHC-13tev-incxsec-1} 把 ATLAS~\cite{Aad:2016naf} 13~TeV 的 $W^+/W^-$ 与 $W/Z$
    fiducial 区比值测量与理论预言比较，含与不含电弱修正两种。
    %    
可以看到无论 $W^+/W^-$ 还是 $W/Z$ 比，所有 PDF 集合都与数据合理符合。
图中理论预言所示不确定度仅为 PDF 不确定度：参数不确定度（$\alpha_s$ 与 $m_c$ 取值）与缺失高阶
QCD 不确定度未计入。
有趣的是，电弱修正使理论预言移动约 0.5\%，并且对所有 PDF 集合都改善了与 ATLAS
测量的一致性。
%
NNPDF3.1 的 PDF 不确定度小于 NNPDF3.0，与 ATLAS 数据符合也更好。

相应的绝对 $W^+$、$W^-$ 与 $Z$ 截面示于图~\ref{fig:LHC-13tev-incxsec-2}，
每种情形都以实验中心值归一。
预言总体与数据相符，可能的例外是 $Z$ 产生的 ABMP16。
与截面比相比，电弱修正对绝对截面的效应——$W^\pm$ 约 1\%、$Z$ 约 0.5\%——
在所涉不确定度的尺度上不太重要，也不必然带来更好的一致性。
    %
把 NNPDF3.1 与 NNPDF3.0 比较再次见到不确定度大幅缩小、预言与数据符合改善。
这一改善对 $Z$ 产生尤为醒目：3.0 当年比数据低约 5\%，
而如今 3.1 在不确定度内相符。
%


%%%%%%%%%%%%%%%%%%%%%%%%%%%%%%%%%%%%%%%%%%%%%%%%%%%%%%%%%%%%%%%%%%%%%
\begin{figure}[t]
\begin{center}
  \includegraphics[scale=1.0]{images/plots/Absolute_ATLAS_shaded.pdf}
  \caption{\small 同图~\ref{fig:LHC-13tev-incxsec-1}，此处为绝对 $W^+$、$W^-$ 与 $Z$ 截面。
    所有预言均以实验中心值归一。
\label{fig:LHC-13tev-incxsec-2}}
\end{center}
\end{figure}
%%%%%%%%%%%%%%%%%%%%%%%%%%%%%%%%%%%%%%%%%%%%%%%%%%%%%%%%%%%%%%%%%%%%%%%

%%%%%%%%%%%%%%%%%%%%%%%%%%%%%%%%%%%%%%%%%%%%%%%%%%%%%%%%%%%%%%%%%%%%%
\begin{figure}[t]
\begin{center}
  \includegraphics[scale=1.0]{images/plots/ggH_ratio.pdf}
  \includegraphics[scale=1.0]{images/plots/ttH_ratio.pdf}
   \includegraphics[scale=1.0]{images/plots/WH_ratio.pdf}
  \includegraphics[scale=1.0]{images/plots/ZH_ratio.pdf}
  \caption{\small LHC $13$ TeV 上 gluo fusion、$t\bar{t}$
    伴随产生与 $VH$ 伴随产生希格斯产生截面的 PDF 依赖。
    所有结果以 NNPDF3.1 中心结果之比给出。
    只显示 PDF 不确定度。
  \label{fig:ggH}
  }
\end{center}
\end{figure}
%%%%%%%%%%%%%%%%%%%%%%%%%%%%%%%%%%%%%%%%%%%%%%%%%%%%%%%%%%%%%%%%%%%%%%%

%%%%%%%%%%%%%%%%%%%%%%%%%%%%%%%%%%%%%%%%%%%%%%%%%%%%%%%%%%%%%%%%%%%%%
\begin{figure}[t]
\begin{center}
  \includegraphics[scale=1.0]{images/plots/vbf_ratio.pdf}
  \includegraphics[scale=1.0]{images/plots/HH_ratio.pdf}
  \caption{\small 同图~\ref{fig:ggH}，左为矢量玻色子融合单希格斯产生，
    右为 gluo fusion 双希格斯产生。
  \label{fig:ggH2}
  }
\end{center}
\end{figure}
%%%%%%%%%%%%%%%%%%%%%%%%%%%%%%%%%%%%%%%%%%%%%%%%%%%%%%%%%%%%%%%%%%%%%%%


\subsection{希格斯产生}

最后我们研究 LHC $13$ TeV 单举希格斯产生预言以及希格斯对产生的 PDF 依赖——
后者在不远的将来也可能进入 LHC 的可达范围~\cite{Bishara:2016kjn,Behr:2015oqq}。
我们研究 gluo fusion 单希格斯、伴随规范玻色子与顶夸克对的伴随产生、矢量玻色子融合，以及 gluo fusion 双希格斯产生。
每种情形我们都给出以 NNPDF3.1 结果归一的预言，且只显示 PDF 不确定度。所有计算（包括 ABMP16）都用 $\alpha_s=0.118$。

设置如下。
对 gluo fusion 用 {\tt ggHiggs}~\cite{Ball:2013bra,Bonvini:2014jma,Bonvini:2016frm}做 N$^3$LO 计算。
重整化与分解标度取 $\mu_F=\mu_R=m_h/2$，计算采用重标度有效理论。
对 $t\bar{t}$ 伴随产生用 {\tt MadGraph5\_aMC@NLO}~\cite{Alwall:2014hca}，
默认分解与重整化标度 $\mu_R=\mu_F=H_T/2$，$H_T$ 为横质量之和。
对电弱规范玻色子伴随产生用 {\tt vh@nnlo} 程序~\cite{Brein:2012ne}以 NNLO 与默认标度设置计算。
对矢量玻色子融合用 {\tt proVBFH}~\cite{Dreyer:2016oyx,Cacciari:2015jma}
以 N$^3$LO 与默认标度设置计算。最后，FCC 100~TeV 的双希格斯产生用
{\tt MadGraph\_aMC@NLO} 计算。

结果示于图~\ref{fig:ggH}--\ref{fig:ggH2}。
对同为胶子 PDF 驱动的 gluo fusion（$x\sim10^{-2}$）与 $t\bar{t}h$（大 $x$），
各 PDF 集合结果在不确定度内一致；
NNPDF3.0 与 NNPDF3.1 也符合良好，新的预言不确定度更小。
规范玻色子伴随产生的结果离散稍大。
NNPDF3.1 预言比 NNPDF3.0 高约 3\%，不确定度缩小一半，因此两个截面勉强在不确定度内相容。
此外在进入 PDF4LHC15 组合的三个 PDF 集合中，
NNPDF3.0 曾给出最小截面，而 NNPDF3.1 如今给出最高者：
$VH$ 产生由 quark-antiquark 亮度驱动，
而 $M_X\simeq 200$ GeV 处 3.0 到 3.1 的增强其实已可在图~\ref{fig:lumi-31-vs-30} 观察到。
对 VBF 我们同样发现 NNPDF3.1 结果比 NNPDF3.0 大约 2\% 且不确定度更小，
与其他 PDF 集合的符合也更好。
最后，对 gluo fusion 双希格斯产生，NNPDF3.1 的中心值略有增大但与 NNPDF3.0 预言一致，
此处所有 PDF 集合同样符合良好。



